\begin{textsample}{NEG Dim 2 – 2021\_c4 – Score 16.00 – 2021\_c4/enade\_2021\_geo\_b7744036e1.txt}  \label{ex:f2_neg_395}
CAPÍTULO 4 PERCEPÇÃO DA PROVA Nas análises feitas neste capítulo, constam as percepções dos concluintes da Área de Geografia ( Licenciatura e Bacharelado ) sobre a prova aplicada no Enade/2021.
Essas percepções foram mensuradas por meio de nove questões que avaliaram desde o grau de dificuldade da prova até o tempo gasto para concluí -la.
As percepções sobre a prova foram relacionadas ao desempenho dos estudantes, à Grande Região de funcionamento do curso, à Categoria Administrativa e ao tipo de Organização Acadêmica da IES.
O questionário de percepção da prova encontra -se ao final dos Anexos VIIIA ( Licenciatura ) e VIIIB ( Bacharelado ), onde está a reprodução do exame.
O desempenho dos estudantes foi classificado em quatro quartos.
Para tanto, esse desempenho foi ordenado de forma ascendente.
O percentil 25, P25, também conhecido como primeiro quartil, é a nota de desempenho que deixa um quarto ( 25 \% ) dos valores observados abaixo e três quartos acima.
Na Figura 1, apresenta -se uma ilustração desse conceito.
O quarto inferior de desempenho é composto pelas notas abaixo do primeiro quartil.
Já o percentil 75, P75, também conhecido como terceiro quartil, é o valor para o qual há três quartos ( 75 \% ) dos dados abaixo e um quarto acima dele.
O quarto superior de desempenho é composto pelas notas iguais ou acima do terceiro quartil.
O percentil 50, P50, também conhecido como mediana, é o valor que divide as notas em dois conjuntos de igual tamanho.
O segundo quarto inclui valores entre o primeiro quartil ( P25 ) e a mediana.
O terceiro quarto contém os valores entre a mediana ( P50 ) e o terceiro quartil ( P75 ).
Vale ressaltar que percentis, quartis e medianas são pontos que, não obrigatoriamente, pertencem ao conjunto original de dados, ao passo que os quartos são subconjuntos dos dados originais.
A seguir, serão apresentados gráficos com resultados selecionados, relativos às nove questões avaliadas por grupos de estudantes.
Nas barras dos gráficos, apresenta -se o percentual de alunos que assinalaram uma das opções ou a soma das porcentagens daqueles que assinalaram duas ( ou três ) delas.
Por exemplo, para as questões 1 e 2, nos gráficos, é apresentada a porcentagem total de participantes que assinalaram as opções ( D ) Difícil e ( E ) Muito difícil.
Em cada barra, foram assinalados, também, os extremos do intervalo de confiança de 95 \% como linhas verticais unidas por uma linha horizontal na forma da letra H maiúscula.
O estimador de um parâmetro com um certo nível de confiança ( e.g. 95,0 \% ) deve conter o parâmetro no intervalo de confiança em 95 \% das vezes.
Na comparação entre os estimadores dos parâmetros de duas classes de uma dada categoria ( e.g.
Norte e Nordeste nas Grandes Regiões, de primeiro e último quarto dentro de desempenho, de IES Públicas e Privadas, ou de Universidades e Faculdades ), associados aos seus respectivos intervalos de confiança, diz -se que não há diferença estatisticamente significativa entre os parâmetros das duas categorias quando há interseção entre os intervalos de confiança, e que há diferença, se os intervalos de confiança são disjuntos ( para maiores informações vide Glossário ).
Nas Tabelas do Anexo IIA ( Licenciatura ) e do Anexo IIB ( Bacharelado ), são apresentados os valores absolutos e a distribuição percentual18 das alternativas válidas das nove questões, segundo o mesmo recorte de desempenho dos alunos e da Grande Região de funcionamento do curso, de Categoria Administrativa e do tipo de Organização Acadêmica da IES. 4.1 GRAU DE DIFICULDADE DA PROVA 4.1.1 COMPONENTE DE FORMAÇÃO GERAL Ao avaliarem “ Qual o grau de dificuldade desta prova na parte de Formação Geral? ” ( Questão 1 ), 39,9 \% do grupo de inscritos e presentes de Licenciatura optaram pelas alternativas Difícil ou Muito difícil.
Entretanto, para mais da metade dos estudantes ( 53,2 \% ), o Componente de Formação Geral da prova foi avaliado com grau de dificuldade Médio ( Gráfico 4.1a, Gráfico 4.2a e, no Anexo IIA, a Tabela II.1a ).
O percentual de estudantes de Licenciatura que consideraram a prova Difícil ou Muito difícil foi maior na região Sul, onde a proporção foi de 46,0 \%, enquanto a de menor proporção foi a Sudeste, com 34,0 \%.
No Gráfico 4.1a, é possível observar que as diferenças entre a região Nordeste e as regiões Sudeste e Sul são estatisticamente significativas, assim como a diferença entre a região Centro-Oeste e a região Sudeste. diferença entre a região Centro-Oeste e a região Sudeste.
Nas Grandes Regiões, a proporção de presentes à prova que avaliaram o Componente de Formação Geral como de grau de dificuldade Médio esteve entre 48,6 \%, na região Sul, e 56,7 \%, na região Sudeste.
Ao avaliarem “ Qual o grau de dificuldade desta prova na parte de Formação Geral? ” ( Questão 1 ), 26,1 \% do grupo de inscritos e presentes de Bacharelado optaram pelas alternativas Difícil ou Muito difícil.
Entretanto, para mais da metade dos estudantes de Bacharelado ( 63,3 \% ), o Componente de Formação Geral da prova foi considerado com grau de dificuldade Médio ( Gráfico 4.1b, Gráfico 4.2b e, no Anexo IIB, a Tabela II.1b ).
O percentual de estudantes de Bacharelado que consideraram a prova como Difícil ou Muito difícil foi maior na região Norte, onde a proporção foi de 36,4 \%, enquanto a de menor incidência foi a Sudeste, com 22,8 \%.
No Gráfico 4.1b, é possível observar que apenas as diferenças entre a região Sudeste e as regiões Sul e Centro-Oeste são estatisticamente significativas.
Nas Grandes Regiões, a proporção de presentes à prova que consideraram o Componente de Formação Geral como sendo de grau de dificuldade Médio esteve entre 54,5 \%, na região Norte e 67,5 \%, na região Nordeste.
O percentual de alunos de Licenciatura que consideraram o componente de Formação Geral da prova Difícil ou Muito difícil foi decrescente em função dos quartos de desempenho : 47,9 \%, no primeiro quarto, e 28,6 \%, no quarto, o grupo de melhor desempenho na prova.
Nos quartos de desempenho intermediários, a proporção de alunos que consideraram a prova Difícil ou Muito difícil foi de 45,4 \%, no segundo quarto, e de 38,0 \%, no terceiro quarto.
As diferenças entre os quartos de desempenho são estatisticamente significativas, com exceção da diferença entre os quartos inferiores de desempenho.
Para todos os quartos de desempenho, a alternativa modal para essa pergunta foi a Médio, com 44,8 \% e 62,6 \% dos respondentes nos quartos extremos, primeiro e quarto, respectivamente.
O percentual de alunos de Bacharelado que consideraram o componente de Formação Geral da prova Difícil ou Muito difícil foi decrescente em função dos quartos de desempenho : 34,0 \%, no primeiro quarto e 20,3 \%, no quarto, grupo de melhor desempenho na prova.
Nos quartos de desempenho intermediários, a proporção de alunos que consideraram a prova Difícil ou Muito difícil foi de 24,2 \%, no segundo quarto e de 25,9 \%, no terceiro quarto.
As diferenças entre o quarto inferior em relação ao segundo quarto e ao último quarto de desempenho são estatisticamente significativas.
Para todos os quartos de desempenho, a alternativa modal para essa pergunta foi a Médio, com 59,3 \% e 66,1 \% dos respondentes no primeiro e no segundo quartos, respectivamente.
O percentual de estudantes de Licenciatura que consideraram o componente de Formação Geral da prova Difícil ou Muito difícil foi maior entre os estudantes de Faculdades, alcançando 53,2 \%, enquanto entre os estudantes dos CEFET/IF, a Organização Acadêmica com a menor incidência, alcançou 38,5 \%.
No Gráfico 4.3a, é possível observar que as diferenças entre Universidades em relação aos Centros Universitários ( 45,4 \% ) e às Faculdades são estatisticamente significativas, assim como entre CEFET/IF em relação aos Centros Universitários e às Faculdades.
Nas Organizações Acadêmicas, a proporção de presentes de Licenciatura à prova que avaliaram o Componente de Formação Geral como de grau de dificuldade Médio esteve entre 41,6 \%, nas Faculdades e 54,2 \% nas Universidades e CEFET/IF ( ver também a Tabela II.2a, no Anexo IIA ).
O percentual de estudantes de Bacharelado que consideraram o componente de Formação Geral da prova como Difícil ou Muito difícil foi maior para os estudantes de Centros Universitários, para os quais a proporção foi de 45,1 \%, enquanto as Universidades foram a Organização Acadêmica com a menor incidência, com 25,3 \%.
No Gráfico 4.3b, é possível observar que as diferenças entre Universidades e Centros Universitários são estatisticamente significativas.
Nas Organizações Acadêmicas, a proporção de presentes de Bacharelado à prova que consideraram o Componente de Formação Geral como sendo de grau de dificuldade Médio foi 63,9 \% nas Universidades e 46,5 \% nos Centros Universitários ( ver também a Tabela II.2b, no Anexo IIB ).
O percentual de estudantes de Licenciatura que consideraram o componente de Formação Geral da prova Difícil ou Muito difícil foi maior entre os estudantes de cursos de IES Privadas ( 46,0 \% ), com diferença estatisticamente significativa da proporção de IES Públicas ( 36,9 \% ).
Já a proporção de presentes de Licenciatura à prova que avaliaram o Componente de Formação Geral como de grau de dificuldade Médio foi de 55,6 \% para os de IES Públicas e de 48,4 \% para os de IES Privadas ( ver Gráfico 4.4a ).
O percentual de estudantes de Bacharelado que consideraram o componente de Formação Geral da prova como Difícil ou Muito difícil foi maior para os estudantes de cursos de IES Privadas ( 34,5 \% ), sem uma diferença estatisticamente significativa da proporção em IES Públicas ( 24,5 \% ).
Já a proporção de presentes de Bacharelado à prova que consideraram o Componente de Formação Geral como sendo de grau de dificuldade Médio foi 63,9 \% para os de IES Públicas e 55,9 \% para os de IES Privadas ( ver Gráfico 4.4b ). 4.1.2 COMPONENTE DE CONHECIMENTO ESPECÍFICO Ao responderem à Questão 2 – “ Qual o grau de dificuldade desta prova na parte de Componente Específico? ” – 44,1 \% do grupo de estudantes de Licenciatura classificaram -na como Difícil ou Muito difícil.
Além disso, o Componente modal de Conhecimento Específico da prova foi considerado de grau de dificuldade Médio por 50,6 \% dos alunos ( Gráfico 4.5a, Gráfico 4.6a, e, no Anexo IIA, a Tabela II.3a ).
A análise das respostas dos estudantes de Licenciatura quanto ao grau de dificuldade do Componente de Conhecimento Específico da prova, agregado por Grande Região, mostra diferenças estatisticamente significativas entre proporções de alunos que a avaliaram como Difícil ou Muito difícil.
Tais diferenças são observadas entre a região com a maior proporção, a Sul ( 50,1 \% ), e as regiões Norte ( 39,9 \% ), Nordeste ( 41,5 \% ), Sudeste ( 41,1 \% ) e Centro-Oeste ( 44,1 \% ).
O percentual de alunos que classificaram o grau de dificuldade como Médio, no Componente de Conhecimento Específico, variou de 45,8 \% a 54,1 \%, nas regiões Sul e Norte, respectivamente.
Ao responderem à Questão 2 – “ Qual o grau de dificuldade desta prova na parte de Componente Específico? ” – 39,2 \% do grupo de estudantes de Bacharelado classificaram -na como Difícil ou Muito difícil.
Além disso, o Componente modal de Conhecimento Específico da prova foi considerado com grau de dificuldade Médio por 54,2 \% dos alunos ( Gráfico 4.5b, Gráfico 4.6b, e, no Anexo IIB, a Tabela II.3b ).
A análise das respostas dos estudantes de Bacharelado quanto ao grau de dificuldade do Componente de Conhecimento Específico da prova, agregado por Grande Região, mostra que a diferença das proporções de alunos que a avaliaram como Difícil ou Muito difícil entre as regiões não é estatisticamente significativa.
O percentual de alunos que classificaram o grau de dificuldade como Difícil ou Muito difícil, no Componente de Conhecimento Específico, variou de 36,4 \% a 45,3 \%, para as regiões Nordeste e Norte, respectivamente.
O percentual de alunos que classificaram o grau de dificuldade como Médio, no Componente de Conhecimento Específico, variou de 45,3 \% a 55,5 \%, para as regiões Norte e Sudeste, respectivamente.
Considerando -se a avaliação da dificuldade das questões do Componente de Conhecimento Específico da prova de acordo com o desempenho dos estudantes de Licenciatura, observa -se uma tendência decrescente dos que consideraram a parte específica Difícil ou Muito difícil, variando de 41,0 \% ( último quarto ) a 47,0 \% ( primeiro quarto ).
As demais proporções foram de 45,6 \% e 42,8 \%, respectivamente, no segundo e terceiro quartos.
Observa -se que há diferença estatisticamente significativa dos resultados entre o quarto inferior de desempenho e os dois quartos superiores de desempenho, assim como entre o segundo quarto e o último quarto.
Já os que responderam que o grau de dificuldade das questões do Componente de Conhecimento Específico da prova foi Médio atingiram percentagens que variaram de 46,1 \%, no primeiro quarto, a 53,8 \%, no último.
Considerando -se a avaliação da dificuldade das questões do Componente de Conhecimento Específico da prova, de acordo com o desempenho dos estudantes de Bacharelado, observa -se que não há diferença estatisticamente significativa dos resultados entre os quartos de desempenho.
A proporção dos que classificaram a parte específica como Difícil ou Muito difícil variou de 36,5 \% ( segundo quarto ) a 42,7 \% ( primeiro quarto ).
As demais proporções dessa reposta foram de 38,3 \% e 39,6 \%, nos quartos terceiro e superior, respectivamente.
Já a proporção dos que responderam que o grau de dificuldade das questões do Componente de Conhecimento Específico da prova foi Médio variou de 51,1 \%, no primeiro quarto a 57,6 \% no segundo.
Considerando -se a avaliação da dificuldade das questões do Componente de Conhecimento Específico da prova, de acordo com o desempenho dos estudantes de Bacharelado, observa -se que não há diferença estatisticamente significativa dos resultados entre os quartos de desempenho.
A proporção dos que classificaram a parte específica como Difícil ou Muito difícil variou de 36,5 \% ( segundo quarto ) a 42,7 \% ( primeiro quarto ).
As demais proporções dessa reposta foram de 38,3 \% e 39,6 \%, nos quartos terceiro e superior, respectivamente.
Já a proporção dos que responderam que o grau de dificuldade das questões do Componente de Conhecimento Específico da prova foi Médio variou de 51,1 \%, no primeiro quarto a 57,6 \% no segundo.
Considerando -se a avaliação da dificuldade das questões do Componente de Conhecimento Específico da prova, o percentual de estudantes de Licenciatura que as consideraram Difícil ou Muito difícil foi maior entre os de Centros Universitários, atingindo 50,9 \% ; já entre os de Universidades, a Organização Acadêmica com a menor incidência, essa proporção foi de 42,5 \%.
No Gráfico 4.7a, é possível observar que a diferença de Universidades em relação às Faculdades ( 50,4 \% ) e aos Centros Universitários é estatisticamente significativa.
Nas Organizações Acadêmicas, a proporção de presentes de Licenciatura à prova que consideraram o Componente de Conhecimento Específico de grau de dificuldade Médio esteve entre 44,8 \%, nas Faculdades, e 51,8 \%, nas Universidades ( ver também Tabela II.4a, no Anexo IIA ).
Considerando -se a avaliação da dificuldade das questões do Componente de Conhecimento Específico da prova, o percentual de estudantes de Bacharelado que consideraram -nas como Difícil ou Muito difícil foi maior para os estudantes de Centros Universitários, para os quais a proporção foi de 54,9 \%, enquanto as Universidades foram a Organização Acadêmica com a menor incidência, com 38,6 \%.
No Gráfico 4.7b, é possível observar que as diferenças entre Universidades e Centros Universitários são estatisticamente significativas.
Nas Organizações Acadêmicas, a proporção de presentes de Bacharelado à prova que consideraram o Componente de Formação Geral como sendo de grau de dificuldade Médio foi 54,7 \% nas Universidades e 42,3 \% nos Centros Universitários ( ver também Tabela II.4b, no Anexo IIB ).
O percentual de estudantes de Licenciatura que consideraram o Componente Específico da prova Difícil ou Muito difícil foi maior entre os estudantes de cursos de IES Privadas ( 50,4 \% ), com diferença estatisticamente significativa da proporção em IES Públicas ( 41,0 \% ).
Já a proporção de presentes de Licenciatura à prova que avaliaram este componente como de grau de dificuldade Médio foi de 53,4 \% para os de IES Públicas e de 45,1 \% para os de IES Privadas ( ver Gráfico 4.8a e Tabela II.4a, no Anexo IIA ).
O percentual de estudantes de Bacharelado que consideraram o Componente Específico da prova como Difícil ou Muito difícil foi maior para os estudantes de cursos de IES Privadas ( 51,0 \% ), uma diferença estatisticamente significativa da proporção em IES Públicas ( 38,3 \% ).
Já a proporção de presentes de Bacharelado à prova que consideraram este componente como sendo de grau de dificuldade Médio foi 55,0 \% para os de IES Públicas e 44,1 \% para os de IES Privadas ( ver Gráfico 4.8b e Tabela II.4b, no Anexo IIB ). 4.2 EXTENSÃO DA PROVA EM RELAÇÃO AO TEMPO TOTAL Indagados quanto à extensão da prova, em relação ao tempo total oferecido para a sua resolução ( Questão 3 ), os estudantes de Licenciatura apontaram, com maior incidência, a alternativa que considerava a extensão Adequada, para todas as agregações consideradas ( Gráfico 4.9a, Gráfico 4.10a, e, no Anexo IIA, a Tabela II.5a ).
O percentual de alunos de Licenciatura que responderam ser a extensão da prova Adequada foi de 49,5 \%.
Já 41,4 \% dos inscritos presentes consideraram que a prova foi Longa ou Muito longa, e 9,2 \% a avaliaram como Curta ou Muito curta.
Entre as Grandes Regiões, os que avaliaram a prova como Longa ou Muito longa, em relação ao tempo total destinado à sua resolução, atingiram uma porcentagem que variou de 31,1 \%, na região Norte, até 42,9 \%, na região Sudeste.
Há diferença estatisticamente significativa entre a região Norte e as regiões Nordeste ( 40,9 \% ), Sudeste e Sul ( 42,8 \% ), e também entre a região Centro-Oeste ( 37,2 \% ) e as regiões Sudeste e Sul.
Indagados quanto à extensão da prova, em relação ao tempo total oferecido para a sua resolução ( Questão 3 ), os estudantes de Bacharelado apontaram, com maior incidência, a alternativa que considerava a extensão Adequada, para todas as agregações consideradas ( Gráfico 4.9b, Gráfico 4.10b, e, no Anexo IIB, a Tabela II.5b ).
O percentual de alunos de Bacharelado que responderam ser a extensão da prova Adequada foi de 50,8 \%.
Já 40,3 \% dos inscritos presentes consideraram que a prova foi Longa ou Muito longa, e 8,9 \% a avaliaram como Curta ou Muito curta.
Entre as Grandes Regiões, a proporção daqueles que avaliaram a prova como Longa ou Muito longa, em relação ao tempo total destinado à sua resolução, variou de 30,9 \%, na região Centro-Oeste até 42,6 \%, na região Sul.
Há diferença estatisticamente significativa entre a região Centro-Oeste e a região Sul.
Considerando -se o desempenho dos alunos de Licenciatura, nota -se ainda que 45,3 \% consideraram a extensão da prova Adequada, no quarto de desempenho inferior, e 52,6 \%, no de melhor desempenho ( quarto superior ).
Nos quartos intermediários, essa proporção foi de 49,4 \%, no segundo quarto, e de 50,5 \%, no terceiro.
No Gráfico 4.10a, observa -se que as proporções de estudantes de Licenciatura que consideraram a prova Longa ou Muito longa apresentam tendência decrescente em função dos quartos de desempenho, sendo a maior proporção dessa resposta no primeiro quarto de desempenho ( 44,0 \% ), e a menor, no quarto superior ( 40,2 \% ).
Pode -se constatar que não há diferença estatisticamente significativa entre as proporções dos quartos de desempenho.
Considerando -se o desempenho dos alunos de Bacharelado, nota -se ainda que 48,2 \% consideraram a extensão da prova Adequada, no segundo quarto de desempenho e 54,7 \%, no de melhor desempenho ( quarto superior ).
No primeiro quarto, essa proporção foi de 48,3 \%, e 51,9 \%, no terceiro quarto.
No Gráfico 4.10b, pode -se constatar que não há diferença estatisticamente significativa entre as proporções de estudantes de Bacharelado que consideraram a prova Longa ou Muito longa em função dos quartos de desempenho, sendo a menor proporção dessa resposta no último quarto de desempenho ( 38,7 \% ), e a maior, no segundo quarto ( 43,2 \% ).
Considerando -se a extensão da prova, o percentual de estudantes de Licenciatura que avaliaram a prova como Longa ou Muito longa, em relação ao tempo total destinado à sua resolução, foi maior entre os estudantes de CEFET/IF, alcançando 45,9 \%, enquanto entre os estudantes de Faculdades, a Organização Acadêmica com a menor incidência, esse índice atingiu 38,8 \%.
No Gráfico 4.11a, é possível observar que não há diferença estatisticamente significativa entre os tipos de Organização Acadêmica.
A proporção de presentes de Licenciatura à prova que consideraram a sua extensão Adequada, em relação ao tempo total destinado à sua resolução, esteve entre 43,0 \%, nos CEFET/IF, e 51,3 \%, nas Faculdades ( ver também Tabela II.6a, no Anexo IIA ).
Considerando -se a extensão da prova, o percentual de estudantes de Bacharelado que avaliaram a prova como Longa ou Muito longa, em relação ao tempo total destinado à sua resolução foi maior para os estudantes de Universidades, para os quais a proporção foi de 40,7 \%, enquanto os Centros Universitários foram a Organização Acadêmica com a menor incidência, com 30,0 \%.
No Gráfico 4.11b, é possível observar que as diferenças entre Universidades e Centros Universitários não são estatisticamente significativas.
Nas Organizações Acadêmicas, a proporção de presentes de Bacharelado à prova que consideraram a sua extensão como Adequada, em relação ao tempo total destinado à sua resolução foi 58,6 \% nos Centros Universitários e 50,5 \% nas Universidades ( ver também Tabela II.6b, no Anexo IIB ).
O percentual de estudantes de Licenciatura que consideraram a prova Longa ou Muito longa, em relação ao tempo total destinado à sua resolução, foi menor entre os estudantes de cursos de IES Privadas ( 41,1 \% ), não havendo diferença estatisticamente significativa da proporção em relação aos de IES Públicas ( 41,5 \% ).
Já a proporção de presentes de Licenciatura à prova que considerou a extensão Adequada foi de 49,0 \% para os de IES Públicas e de 50,4 \% para os de IES Privadas ( ver Gráfico 4.12a e Tabela II.6a, no Anexo IIA ).
O percentual de estudantes de Bacharelado que consideraram a prova como Longa ou Muito longa, em relação ao tempo total destinado à sua resolução foi menor para os estudantes de cursos de IES Privadas ( 33,3 \% ), sem diferença estatisticamente significativa da proporção em IES Públicas ( 40,8 \% ).
Já a proporção de presentes de Bacharelado à prova que consideraram a extensão como Adequada foi 50,4 \% para os de IES Públicas e 54,9 \% para os de IES Privadas ( ver Gráfico 4.12b e Tabela II.6b, no Anexo IIB ). 4.3 COMPREENSÃO DOS ENUNCIADOS DAS QUESTÕES 4.3.1 COMPONENTE DE FORMAÇÃO GERAL Em relação aos enunciados das questões do Componente de Formação Geral ( Questão 4 ), as opiniões foram positivas, já que 71,8 \% dos alunos de Licenciatura avaliados consideraram todos ou a maioria dos enunciados das questões “ claros e objetivos ” ( Gráfico 4.13a, Gráfico 4.14a e, no Anexo IIA, a Tabela II.7a ).
Na análise regional, a percentagem de estudantes de Licenciatura que avaliaram que todos ou a maioria dos enunciados das questões do Componente de Formação Geral estavam “ claros e objetivos ” variou de 65,7 \%, na região Norte, a 73,7 \%, na região Sudeste.
A análise das percepções dos estudantes de Licenciatura sobre a clareza e a objetividade dos enunciados permite afirmar que todos, ou a maioria dos enunciados de questões relativas ao Componente de Formação Geral, foram considerados “ claros e Em relação aos enunciados das questões do Componente de Formação Geral ( Questão 4 ), as opiniões foram positivas, já que 75,2 \% dos alunos de Bacharelado avaliados consideraram todos ou a maioria dos enunciados das questões “ claros e objetivos ” ( Gráfico 4.13b, Gráfico 4.14b, e, no Anexo IIB, a Tabela II.7b ).
Na análise regional, a percentagem de estudantes de Bacharelado que avaliaram que todos ou a maioria dos enunciados das questões do Componente de Formação Geral estavam “ claros e objetivos ” variou de 70,1 \%, na região Norte a 76,1 \%, na região Sudeste.
A análise das percepções dos estudantes de Bacharelado sobre a clareza e a objetividade dos enunciados permite afirmar que todos, ou a maioria dos enunciados de questões relativas ao Componente de Formação Geral, foram considerados “ claros e objetivos ” para a maior parte dos respondentes ( maior ou igual a 70,1 \%, em todas as regiões e maior ou igual a 67,4 \% para todos os quartos de desempenho ).
Segundo o desempenho, observa -se que a proporção dos que emitiram essa opinião cresce, conforme o desempenho aumenta, com diferenças estatisticamente significativas entre os quartos de desempenho.
No quarto superior, a clareza e a objetividade de todos ou da maioria dos enunciados das questões foram percebidas por 80,7 \%.
Segundo o desempenho, observa -se que a proporção dos que emitiram essa opinião cresce conforme o desempenho aumenta, com diferenças estatisticamente significativas entre o quarto inferior e os quartos segundo e superior de desempenho.
No quarto superior, a clareza e a objetividade de todos ou da maioria dos enunciados das questões foram percebidas por 80,1 \%.
A proporção dos presentes de Licenciatura que consideraram que todos ou a maioria dos enunciados das questões do Componente de Formação Geral estavam “ claros e objetivos ” foi maior entre os estudantes de Centros Universitários, chegando a 72,3 \%, enquanto entre os de Faculdades, a Organização Acadêmica com a menor incidência, esse índice atingiu 70,3 \%.
No Gráfico 4.15a, é possível observar que as diferenças entre os tipos de Organização Acadêmica não são estatisticamente significativas ( ver também Tabela II.8a, no Anexo IIA ).
A proporção dos presentes de Bacharelado que consideraram que todos ou a maioria dos enunciados das questões do Componente de Formação Geral estavam “ claros e objetivos ” foi maior para os estudantes de Centros Universitários, para os quais a proporção foi de 81,7 \%, enquanto os de Universidades foram a Organização Acadêmica com a menor incidência, com 75,0 \%.
No Gráfico 4.15b, é possível observar que não há diferença estatisticamente significativa entre as Organizações Acadêmicas ( ver também Tabela II.8b, no Anexo IIB ).
O percentual de estudantes de Licenciatura que consideraram que todos ou a maioria dos enunciados das questões do Componente de Formação Geral estavam “ claros e objetivos ” foi menor entre os estudantes de cursos de IES Públicas ( 71,6 \% ), mas sem diferença estatisticamente significativa da proporção em IES Privadas ( 72,3 \% ).
Um maior detalhamento pode ser encontrado no Gráfico 4.16a e na Tabela II.8a, no Anexo IIA.
O percentual de estudantes de Bacharelado que consideraram que todos ou a maioria dos enunciados das questões do Componente de Formação Geral estavam “ claros e objetivos ” foi maior para os estudantes de cursos de IES Privadas ( 80,6 \% ), sem diferença estatisticamente significativa da proporção em IES Públicas ( 74,8 \% ).
Um maior detalhamento pode ser encontrado no Gráfico 4.16b e na Tabela II.8b, no Anexo IIB. 4.3.2 COMPONENTE DE CONHECIMENTO ESPECÍFICO Em relação aos enunciados das questões do Componente de Conhecimento Específico da prova, para 72,8 \% dos estudantes de Licenciatura avaliados da Área de Geografia, a clareza e a objetividade ( Questão 5 ) estavam presentes em todas ou na maioria das questões ( Gráfico 4.17a a Gráfico 4.20a, e no Anexo IIA, as Tabelas II.9a e II.10a ).
A maioria dos estudantes de Licenciatura de todas as Grandes Regiões brasileiras considerou “ claros e objetivos ” todos ou a maioria dos enunciados das questões do Componente de Conhecimento Específico da prova, percentual sempre maior ou igual a 67,8 \%.
A diferença entre a região Sul ( 74,8 \% ) e as regiões Norte ( 67,8 \% ) e Nordeste ( 70,8 \% ) é estatisticamente significativa.
Em relação aos enunciados das questões do Componente de Conhecimento Específico da prova, para 76,3 \% dos estudantes de Bacharelado avaliados da Área de Geografia, a clareza e a objetividade ( Questão 5 ) estavam presentes em todas ou na maioria das questões ( Gráfico 4.17b a Gráfico 4.20b, e no Anexo IIB, as Tabelas II.9b e II.10b ).
A maioria dos estudantes de Bacharelado de todas as Grandes Regiões brasileiras considerou “ claros e objetivos ” todos ou a maioria dos enunciados das questões do Componente de Conhecimento Específico da prova, percentual sempre maior ou igual a 72,3 \%.
A diferença entre as regiões não é estatisticamente significativa.
A proporção de estudantes de Licenciatura que consideraram os enunciados das questões “ claros e objetivos ” apresenta uma tendência crescente em relação ao aumento de desempenho : mais elevada no quarto superior ( 80,4 \% ), se comparada ao quarto inferior de desempenho ( 63,3 \% ).
As diferenças entre todos os quartos de desempenho são estatisticamente significativas.
A proporção de estudantes de Bacharelado que consideraram os enunciados das questões “ claros e objetivos ” apresenta uma tendência crescente em relação ao aumento de desempenho : mais elevada no quarto superior ( 81,1 \% ), se comparada ao quarto inferior de desempenho ( 69,7 \% ).
A diferença entre os quartos extremos de desempenho, o primeiro e o último, é estatisticamente significativa.
A proporção dos presentes de Licenciatura que consideraram que todos ou a maioria dos enunciados das questões do Componente de Conhecimento Específico estavam “ claros e objetivos ” foi maior entre os estudantes de Centros Universitários, atingindo 73,6 \%.
Entre os estudantes de CEFET/IF, a Organização Acadêmica com a menor incidência, essa proporção atingiu 71,1 \%, sendo intermediária a proporção de Universidades ( 72,8 \% ) e Faculdades ( 72,0 \% ).
No Gráfico 4.19a, é possível observar que as diferenças entre os tipos de Organização Acadêmica não são estatisticamente significativas ( ver também Tabela II.10a, no Anexo IIA ).
A proporção dos presentes de Bacharelado que consideraram que todos ou a maioria dos enunciados das questões do Componente de Conhecimento Específico estavam “ claros e objetivos ” foi maior para os estudantes de Centros Universitários, para os quais a proporção foi de 83,1 \%, enquanto as Universidades foram a Organização Acadêmica com a menor incidência, com 76,0 \%.
No Gráfico 4.19b, é possível observar que não há diferença estatisticamente significativa entre Universidades e Centros Universitários ( ver também Tabela II.10b, no Anexo IIB ).
O percentual de estudantes de Licenciatura que consideraram que todos ou a maioria dos enunciados das questões do Componente de Conhecimento Específico estavam “ claros e objetivos ” foi maior entre os estudantes de cursos de IES Públicas ( 72,8 \% ), sem diferença estatisticamente significativa da proporção em IES Privadas ( 72,7 \% ).
Ver Gráfico 4.20a e Tabela II.10a, no Anexo IIA, para maior detalhamento.
O percentual de estudantes de Bacharelado que consideraram que todos ou a maioria dos enunciados das questões do Componente de Conhecimento Específico estavam “ claros e objetivos ” foi maior para os estudantes de cursos de IES Privadas ( 82,1 \% ), sem uma diferença estatisticamente significativa da proporção em IES Públicas ( 75,8 \% ).
Ver Gráfico 4.20b e Tabela II.10b, no Anexo IIB, para um maior detalhamento. 4.4 SUFICIÊNCIA DAS INFORMAÇÕES/INSTRUÇÕES FORNECIDAS Ao se avaliarem as informações/instruções fornecidas para a resolução das questões ( Questão 6 ), constatou -se que 80,8 \% dos respondentes de Licenciatura da Área de Geografia de todo o Brasil afirmaram que estas eram até excessivas ou suficientes em todas ou na maioria das questões ( Gráfico 4.21a, Gráfico 4.22a e, no Anexo IIA, a Tabela II.11a ).
Quanto à distribuição de respondentes pelas Grandes Regiões, observa -se que a proporção de estudantes que consideraram as informações/instruções fornecidas até excessivas ou suficientes em todas ou na maioria das questões foi sempre superior ou igual a 78,1 \%, chegando a 82,4 \%, na região Sudeste.
Existe diferença estatisticamente significativa apenas entre as regiões Sudeste e Nordeste.
Ao avaliarem as informações/instruções fornecidas para a resolução das questões ( Questão 6 ), 85,2 \% dos respondentes de Bacharelado da Área de Geografia de todo o Brasil afirmaram que estas eram até excessivas ou suficientes em todas ou na maioria das questões ( Gráfico 4.21b, Gráfico 4.22b, e, no Anexo IIB, a Tabela II.11b ).
Quanto à distribuição de respondentes pelas Grandes Regiões, observa -se que a proporção de estudantes que consideraram as informações/instruções fornecidas até excessivas ou suficientes em todas ou na maioria das questões foi sempre superior ou igual a 77,9 \%, chegando a 86,9 \%, na região Sul.
Não existem diferenças estatisticamente significativas entre as regiões.
Levando -se em conta o desempenho dos participantes de Licenciatura, observa -se uma tendência crescente das proporções de participantes que avaliaram as informações/instruções como até excessivas ou suficientes em todas ou na maioria das questões, havendo diferenças estatisticamente significativas entre os quartos de desempenho, conforme mostrado no Gráfico 4.22a.
O percentual foi mais elevado no quarto superior ( 86,7 \% ), percentual superior à média nacional ( 80,8 \% ).
Já no quarto inferior, a suficiência das informações/instruções declarada como até excessiva, em todas ou na maioria das questões foi percebida por 72,2 \% dos respondentes.
Levando -se em conta o desempenho dos participantes de Bacharelado, observa -se uma tendência crescente das proporções de participantes que avaliaram as informações/instruções como até excessivas ou suficientes em todas ou na maioria das questões, havendo diferenças estatisticamente significativas entre os quartos extremos de desempenho, o primeiro e o último, como mostrado no Gráfico 4.22b.
O percentual foi mais elevado no quarto superior ( 89,7 \% ), percentual superior à média nacional ( 85,2 \% ).
Já no quarto inferior, a suficiência das informações/instruções declarada como até excessiva, em todas ou na maioria das questões foi percebida por 80,2 \% dos respondentes.
Considerando -se as informações/instruções fornecidas para a resolução das questões ( Questão 6 ), a proporção que considerou que estas eram até excessivas ou suficientes em todas ou na maioria das questões foi maior entre os estudantes de Licenciatura de Universidades, totalizando 81,2 \%.
Os CEFET/IF foram a Organização Acadêmica em que houve a menor proporção, 78,8 \%.
No Gráfico 4.23a, é possível observar que as diferenças entre as Organizações Acadêmicas não são estatisticamente significativas ( ver também Tabela II.12a, no Anexo IIA ).
Considerando -se as informações/instruções fornecidas para a resolução das questões ( Questão 6 ), a proporção que considerou que estas eram até excessivas ou suficientes em todas ou na maioria das questões foi maior para os estudantes de Bacharelado de Universidades, para os quais a proporção foi de 85,2 \%, enquanto os de Centros Universitários foram a Organização Acadêmica com a menor incidência, com 84,5 \%.
No Gráfico 4.23b, é possível observar que não existe diferença estatisticamente significativa entre Universidades e Centros Universitários ( ver também Tabela II.12b, no Anexo IIB ).
O percentual de estudantes de Licenciatura que consideraram que as informações/instruções fornecidas para a resolução das questões eram até excessivas ou suficientes em todas ou na maioria das questões foi maior para os estudantes de cursos de IES Públicas ( 81,0 \% ), não havendo diferença estatisticamente significativa da proporção em IES Privadas ( 80,4 \% ) ( ver Gráfico 4.24a e Tabela II.12a, no Anexo IIA ). 4.5 DIFICULDADE ENCONTRADA AO RESPONDER À PROVA Indagados sobre as dificuldades com as quais se depararam ao responder à prova ( Questão 7 ), 11,0 \% dos estudantes de Licenciatura apontaram o Desconhecimento do conteúdo.
Para 46,9 \%, a Forma diferente de abordagem do conteúdo foi indicada como dificuldade.
Já a Falta de motivação para fazer a prova foi a dificuldade apontada por 20,2 \% dos respondentes.
Considerando -se todo o Brasil, 15,6 \% dos respondentes afirmaram que não tiveram qualquer tipo de dificuldade para responder à prova ( Tabelas II.13a e II.14a, no Anexo IIA ).
Nos Gráficos 4.25a a 4.28a, são apresentados os percentuais de estudantes de Licenciatura que apontaram o Desconhecimento do conteúdo como “ dificuldade ao responder à prova ”.
Na análise por Grandes Regiões, o percentual de inscritos e presentes de Licenciatura que apontaram o Desconhecimento do conteúdo como “ dificuldade ao responder à prova ” não 4.5 DIFICULDADE ENCONTRADA AO RESPONDER À PROVA Indagados sobre as dificuldades com as quais se depararam ao responder à prova ( Questão 7 ), 11,0 \% dos estudantes de Licenciatura apontaram o Desconhecimento do conteúdo.
Para 46,9 \%, a Forma diferente de abordagem do conteúdo foi indicada como dificuldade.
Já a Falta de motivação para fazer a prova foi a dificuldade apontada por 20,2 \% dos respondentes.
Considerando -se todo o Brasil, 15,6 \% dos respondentes afirmaram que não tiveram qualquer tipo de dificuldade para responder à prova ( Tabelas II.13a e II.14a, no Anexo IIA ).
Nos Gráficos 4.25a a 4.28a, são apresentados os percentuais de estudantes de Licenciatura que apontaram o Desconhecimento do conteúdo como “ dificuldade ao responder à prova ”.
Na análise por Grandes Regiões, o percentual de inscritos e presentes de Licenciatura que apontaram o Desconhecimento do conteúdo como “ dificuldade ao responder à prova ” não superou 12,7 \%.
Os percentuais variaram de 10,2 \%, nas regiões Sudeste e Centro-Oeste, a 12,7 \%, na região Sul.
A Forma diferente de abordagem do conteúdo foi a escolha modal dos estudantes de Licenciatura, com percentuais que variaram de 43,0 \% ( região Sudeste ) a 52,2 \% ( região Centro-Oeste ).
O percentual de alunos que citaram a Falta de motivação como dificuldade variou de 15,4 \% ( região Norte ) a 24,0 \% ( região Sudeste ).
Os que declararam não ter qualquer tipo de dificuldade para responder à prova variaram de 10,7 \%, na região Norte, a 17,5 \%, na região Sul.
Indagados sobre as dificuldades com as quais se depararam ao responder à prova ( Questão 7 ), 10,8 \% dos estudantes de Bacharelado apontaram o Desconhecimento do conteúdo.
Para 40,8 \%, a Forma diferente de abordagem do conteúdo foi indicada como dificuldade.
Já a Falta de motivação para fazer a prova foi a dificuldade apontada por 28,9 \% dos respondentes.
Considerando -se todo o Brasil, 13,9 \% dos respondentes afirmaram que não tiveram qualquer tipo de dificuldade para responder à prova ( Tabelas II.13b e II.14b, no Anexo IIB ).
Nos Gráficos 4.25b a 4.28b, são apresentados os percentuais de estudantes de Bacharelado que apontaram o Desconhecimento do conteúdo como “ dificuldade ao responder à prova ”.
Na análise por Grandes Regiões, o percentual de inscritos e presentes de Bacharelado que apontaram o Desconhecimento do conteúdo como “ dificuldade ao responder à prova ” não superou 11,8 \%.
Os percentuais variaram de 9,8 \%, na região Centro-Oeste a 11,8 \%, na região Norte.
A Forma diferente de abordagem do conteúdo foi a escolha modal dos estudantes de Bacharelado, com percentuais que variaram de 37,4 \% ( região Sudeste ) a 53,9 \% ( região Norte ).
O percentual de alunos que citaram a Falta de motivação como dificuldade variou de 19,7 \% ( região Norte ) a 30,3 \% ( regiões Nordeste e Sudeste ).
Os que declararam não ter qualquer tipo de dificuldade para responder à prova variaram de 7,9 \%, na região Norte a 16,7 \%, na região Centro-Oeste.
Em relação aos quartos de desempenho, o Desconhecimento do conteúdo foi a opção escolhida por 13,0 \% dos estudantes de Licenciatura do quarto inferior e por 10,1 \% do quarto superior, caracterizando uma tendência decrescente.
São observadas diferenças estatisticamente significativas entre o quarto inferior e os dois quartos superiores de desempenho.
A alternativa modal para os alunos de Licenciatura, quando agregados pelos quartos de desempenho, para a dificuldade encontrada, foi causada pela Forma diferente de abordagem do conteúdo : 43,9 \%, no quarto superior, e 49,1 \%, no segundo quarto, escolheram essa alternativa.
Em relação aos quartos de desempenho, o Desconhecimento do conteúdo foi a opção escolhida por 9,7 \% dos estudantes de Bacharelado do quarto inferior e por 12,2 \% do quarto superior.
É observada diferença estatisticamente significativa apenas entre os quartos intermediários de desempenho.
A alternativa modal para os alunos de Bacharelado, quando agregados pelos quartos de desempenho, para a dificuldade encontrada, foi causada pela Forma diferente de abordagem do conteúdo : variando de 36,5 \%, no quarto superior, a 44,0 \%, no terceiro quarto.
Na análise por Organização Acadêmica, o percentual de inscritos e presentes de Licenciatura que apontaram o Desconhecimento do conteúdo como “ dificuldade ao responder à prova ” não superou 17,7 \%.
Os percentuais variaram de 7,7 \%, nos CEFET/IF, a 17,7 \%, nas Faculdades.
A Forma diferente de abordagem do conteúdo foi a escolha modal dos estudantes de Licenciatura, com percentuais que variaram de 46,5 \% ( Universidades ) a 49,1 \% ( CEFET/IF ).
O percentual de alunos de Licenciatura que citaram a Falta de motivação como dificuldade variou de 16,3 \% ( Centros Universitários ) a 21,2 \% ( CEFET/IF ).
Os que declararam não ter qualquer tipo de dificuldade para responder à prova variaram de 12,1 \%, nas Faculdades, a 16,2 \%, nos Centros Universitários.
No Gráfico 4.27a, é possível observar que há diferenças estatisticamente significativas de Universidades em relação às demais organizações, assim como de CEFET/IF em relação aos Centros Universitários e às Faculdades ( ver também Tabela II.14a, no Anexo IIA ).
Na análise por Organização Acadêmica, o percentual de inscritos e presentes de Bacharelado que apontaram o Desconhecimento do conteúdo como “ dificuldade ao responder à prova ” não superou 14,1 \%.
Os percentuais foram de 10,6 \% nas Universidades e 14,1 \% nos Centros Universitários.
A Forma diferente de abordagem do conteúdo foi a escolha modal dos estudantes de Bacharelado, com percentuais que foram de 40,7 \% ( Universidades ) e 45,1 \% ( Centros Universitários ).
O percentual de alunos de Bacharelado que citaram a Falta de motivação como dificuldade foi de 9,9 \% ( Centros Universitários ) e 29,6 \% ( Universidades ).
Os que declararam não ter qualquer tipo de dificuldade para responder à prova foram de 13,5 \%, nas Universidades e 23,9 \%, nos Centros Universitários.
No Gráfico 4.27b, é possível observar que a diferença entre Universidades e Centros Universitários não é estatisticamente significativa ( ver também Tabela II.14b, no Anexo IIB ).
Em relação às Categorias Administrativas, o Desconhecimento do conteúdo foi a opção escolhida por 10,4 \% dos estudantes de Licenciatura nas IES Públicas e por 12,4 \% nas IES Privadas.
São observadas diferenças estatisticamente significativas entre essas duas categorias.
A alternativa modal para os alunos de Licenciatura, quando agregados pelas Categorias Administrativas, para a dificuldade encontrada, foi causada pela Forma diferente de abordagem do conteúdo : 46,3 \% para IES Públicas e 48,2 \% para as IES Privadas ( ver Gráfico 4.28a e Tabela II.14a, no Anexo IIA ).
Em relação às Categorias Administrativas, o Desconhecimento do conteúdo foi a opção escolhida por 10,6 \% dos estudantes de Bacharelado nas IES Públicas e por 13,1 \% nas IES Privadas.
Não são observadas diferenças estatisticamente significativas entre essas duas categorias.
A alternativa modal para os alunos de Bacharelado, quando agregados pelas Categorias Administrativas, para a dificuldade encontrada, foi causada pela Forma diferente de abordagem do conteúdo : 40,4 \% para IES Públicas e 46,2 \% para as IES Privadas ( ver Gráfico 4.28b e Tabela II.14b, no Anexo IIB ). 4.6 CONTEÚDOS DAS QUESTÕES OBJETIVAS DA PROVA Ao se analisarem os conteúdos das questões objetivas da prova ( Questão 8 ), verifica -se que um percentual pequeno dos estudantes de Licenciatura avaliados, apenas 6,2 \%, afirmou que não estudou ainda a maioria desses conteúdos ( Gráficos 4.29a a Gráfico 4.32a, e as Tabelas II.15a e II.16a, no Anexo IIA ).
A maioria ( 64,2 \% ) afirmou ter estudado e aprendido muitos ou todos os conteúdos avaliados.
Na análise por Grande Região, a proporção de respondentes de Licenciatura que escolheram a opção não estudou ainda a maioria desses conteúdos foi pequena.
Observa -se que, nas regiões Norte ( 7,0 \% ) e Sul ( 8,5 \% ), as proporções foram maiores do que a média nacional ( 6,2 \% ).
Nas regiões Nordeste, Sudeste e Centro-Oeste, a situação foi inversa, com proporções mais baixas, correspondendo a, respectivamente, 5,1 \%, 5,4 \% e 4,9 \%.
São observadas diferenças estatisticamente significativas entre a região Sul e as regiões Nordeste, Sudeste e Centro-Oeste.
Em todas as regiões, a maioria dos presentes afirmou ter estudado e aprendido muitos ou todos os conteúdos, com proporções variando entre 57,2 \%, na região Norte, a 70,5 \%, na região Sudeste.
Ao analisar os conteúdos das questões objetivas da prova ( Questão 8 ), um percentual muito pequeno dos estudantes de Bacharelado avaliados, apenas 2,5 \%, afirmou que não estudou ainda a maioria desses conteúdos ( Gráficos 4.29b a Gráfico 4.32b, e as Tabelas II.15b e II.16b, no Anexo IIB ).
A maioria ( 74,0 \% ) afirmou ter estudado e aprendido muitos ou todos os conteúdos avaliados.
Na análise por Grande Região, a proporção de respondentes de Bacharelado que escolheram a opção não estudou ainda a maioria desses conteúdos foi pequena.
Observa -se que, nas regiões Norte ( 6,5 \% ) e Centro-Oeste ( 3,4 \% ), as proporções foram maiores do que a média nacional ( 2,5 \% ).
Nas regiões Nordeste ( 1,7 \% ), Sudeste ( 2,4 \% ) e Sul ( 2,3 \% ), a situação foi inversa : uma proporção mais baixa.
Não são observadas diferenças estatisticamente significativas entre as regiões.
Em todas as regiões, a maioria dos presentes afirmou ter estudado e aprendido muitos ou todos os conteúdos, com proporções variando entre 61,0 \%, na região Norte a 77,1 \%, na região Sudeste.
Considerando -se separadamente as opiniões de estudantes de Licenciatura dos quatro quartos de desempenho, observa -se que, no quarto inferior, 11,8 \% ofereceram como resposta que não estudou ainda a maioria desses conteúdos, sendo 1,6 \% os do quarto superior com a mesma resposta.
A diferença entre os alunos que optaram por esse motivo de dificuldade nos quartos de desempenho é estatisticamente significativa.
Tendo -se em conta o quarto superior, 81,2 \% dos alunos de Licenciatura afirmaram ter estudado e aprendido muitos ou todos os conteúdos.
No outro extremo, no primeiro quarto, 45,7 \% optaram pelas mesmas categorias.
As proporções são crescentes com o desempenho.
Considerando -se, separadamente, as opiniões de estudantes de Bacharelado dos quatro quartos de desempenho, observa -se que, no quarto inferior, 5,6 \% ofereceram como resposta que não estudou ainda a maioria desses conteúdos, sendo 0,6 \% os do quarto superior com a mesma resposta.
A diferença entre os alunos que optaram por esse motivo de dificuldade no quarto inferior, em relação aos demais quartos de desempenho, é estatisticamente significativa.
Tendo -se em conta o quarto superior, 85,4 \% dos alunos de Bacharelado afirmaram ter estudado e aprendido muitos ou todos os conteúdos.
No outro extremo, no primeiro quarto, 61,2 \% optaram pelas mesmas categorias.
As proporções são crescentes com o desempenho.
Na análise por Organização Acadêmica, a proporção de respondentes de Licenciatura que escolheu a opção não estudou ainda a maioria desses conteúdos foi pequena.
Observa -se que as proporções nos Centros Universitários e Faculdades, com respectivamente 12,4 \% e 12,5 \%, foram maiores do que a média nacional ( 6,2 \% ).
Nas Universidades e CEFET/IF, a situação foi inversa : a proporção foi mais baixa ( 5,0 \% e 4,1 \%, respectivamente ).
No Gráfico 4.31a, é possível observar que as diferenças de Universidades em relação a Centros Universitários e Faculdades são estatisticamente significativas, assim como de CEFET/IF em relação a Centros Universitários e Faculdades.
Em todos os tipos de Organização Acadêmica, a maioria dos presentes afirmou ter estudado e aprendido muitos ou todos os conteúdos, com proporções variando entre 46,1 \%, nas Faculdades, a 69,3 \%, nos CEFET/IF ( ver também Tabela II.16a, no Anexo IIA ).
Na análise por Organização Acadêmica, a proporção de respondentes de Bacharelado que escolheram a opção não estudou ainda a maioria desses conteúdos foi pequena.
Observa -se que, nas Universidades ( 2,4 \% ) a proporção foi menor do que a média nacional ( 2,5 \% ).
Nos Centros Universitários ( 2,8 \% ), a situação foi inversa : a proporção foi mais alta.
No Gráfico 4.31b, é possível observar que a diferença entre Universidades e Centros Universitários não é estatisticamente significativa ( ver também Tabela II.16b, no Anexo IIB ).
Em todos os tipos de Organização Acadêmica, a maioria dos presentes afirmou ter estudado e aprendido muitos ou todos os conteúdos, com proporções de 73,8 \%, nas Universidades e de 77,5 \%, nos Centros Universitários.
Em relação às Categorias Administrativas, observa -se que, nas IES Públicas, 4,5 \% ofereceram como resposta que não estudou ainda a maioria desses conteúdos, sendo 9,8 \% nas IES Privadas com a mesma resposta.
A diferença entre os alunos de Licenciatura que optaram por esse motivo de dificuldade nas duas categorias é estatisticamente significativa.
Considerando -se as IES Públicas, 67,5 \% dos alunos de Licenciatura afirmaram ter estudado e aprendido muitos ou todos os conteúdos, sendo de 57,7 \% a proporção dos que escolheram essa resposta nas IES Privadas ( ver Gráfico 4.32a e Tabela II.16a, no Anexo IIA ).
Em relação às Categorias Administrativas, observa -se que, nas IES Públicas, 2,4 \% ofereceram como resposta que não estudou ainda a maioria desses conteúdos, sendo 3,5 \% nas IES Privadas com a mesma resposta.
A diferença entre os alunos de Bacharelado que optaram por esse motivo de dificuldade nas duas categorias não é estatisticamente significativa.
Considerando -se as IES Públicas, 73,4 \% dos alunos de Bacharelado afirmaram ter estudado e aprendido muitos ou todos os conteúdos.
Na outra categoria, 81,3 \% optaram pelas mesmas categorias ( ver Gráfico 4.32b e Tabela II.16b, no Anexo IIB ). 4.7 TEMPO GASTO PARA CONCLUIR A PROVA Ao responderem sobre o tempo de conclusão da prova ( Questão 9 ), mais de três quartos dos estudantes de Licenciatura ( 76,1 \% ) afirmaram ter gasto Entre duas e quatro horas ( Gráfico 4.33a a Gráfico 4.36a e, no Anexo IIA, as Tabelas II.17a e II.18a ).
Considerando -se as cinco Grandes Regiões brasileiras, os percentuais dos que utilizaram Entre duas e quatro horas para finalizar a prova nas regiões Sudeste ( 75,5 \% ) e Sul ( 75,5 \% ) foram inferiores ao percentual nacional.
Nas demais Grandes Regiões, os percentuais de alunos de Licenciatura que dispensaram Entre duas e quatro horas para concluir a prova foram : 77,5 \%, na região Norte, 76,9 \%, na região Nordeste e 77,0 \%, na região Centro-Oeste, como mostrado no Gráfico 4.33a.
Não há diferenças estatisticamente significativas entre as regiões.
Ao responderem sobre o tempo de conclusão da prova ( Questão 9 ), quase três quartos dos estudantes de Bacharelado ( 74,9 \% ) afirmaram ter gasto Entre duas e quatro horas ( Gráfico 4.33b a Gráfico 4.36b e, no Anexo IIB, as Tabelas II.17b e II.18b ).
Considerando -se as cinco Grandes Regiões brasileiras, os percentuais dos que utilizaram Entre duas e quatro horas para finalizar a prova nas regiões Sudeste ( 72,9 \% ) e Sul ( 70,7 \% ) foram inferiores ao percentual nacional.
Nas demais Grandes Regiões, o percentual de alunos de Bacharelado que dispensaram Entre duas e quatro horas para concluir a prova foi : 84,4 \%, na região Norte, 78,8 \%, na região Nordeste e 78,9 \%, na região Centro-Oeste, como mostrado no Gráfico 4.33b.
A diferença entre as regiões Nordeste e Sul é estatisticamente significativa, assim como entre a região Norte e as regiões Sudeste e Sul.
Uma vez analisadas as alternativas escolhidas pelos estudantes de Licenciatura que se situam nos diferentes quartos de desempenho, observa -se uma tendência crescente : uma proporção maior de participantes no quarto superior declarou ter gasto Entre duas e quatro horas para concluir a prova quando comparada com a proporção de participantes nos quartos inferiores.
Há diferenças estatisticamente significativas entre o quarto inferior e os demais quartos de desempenho, bem como entre o segundo e o quarto superior de desempenho.
Uma vez analisadas as alternativas escolhidas pelos estudantes de Bacharelado que se situam nos diferentes quartos de desempenho, observa -se uma tendência crescente : uma proporção maior de participantes no quarto superior declarou ter gasto Entre duas e quatro horas para concluir a prova quando comparada com a proporção de participantes nos quartos inferiores.
Há diferenças estatisticamente significativas entre quarto inferior e os demais quartos de desempenho.
Na análise por Organização Acadêmica, o percentual dos que utilizaram Entre duas e quatro horas para finalizar a prova nas Universidades ( 76,4 \% ) foi superior ao percentual nacional ( 76,1 \% ).
Nos Centros Universitários ( 75,1 \% ), Faculdades ( 74,9 \% ) e CEFET/IF ( 75,4 \% ), tal percentual foi inferior, e, como mostrado no Gráfico 4.35a, não há diferença estatisticamente significativa entre as proporções dos tipos de Organização Acadêmica.
Na análise por Organização Acadêmica, os percentuais dos que utilizaram Entre duas e quatro horas para finalizar a prova nas Universidades e Centros Universitários foram : 75,0 \% e 72,5 \%, respectivamente.
Como mostrado no Gráfico 4.35b, a diferença entre Universidades e Centros Universitários não é estatisticamente significativa.
Em relação às Categorias Administrativas, observa -se que, nas IES Públicas, 76,2 \% declararam ter gasto Entre duas e quatro horas para concluir a prova, sendo 76,0 \% nas IES Privadas com a mesma resposta.
A diferença entre os alunos de Licenciatura que optaram por esse motivo de dificuldade nas duas categorias não é estatisticamente significativa ( ver Gráfico 4.36a e Tabela II.18a, no Anexo IIA ).
Em relação às Categorias Administrativas, observa -se que, nas IES Públicas, 74,9 \% declararam ter gasto Entre duas e quatro horas para concluir a prova, sendo 75,5 \% nas IES Privadas com a mesma resposta.
A diferença entre os alunos de Bacharelado que optaram por esse motivo de dificuldade nas duas categorias não é estatisticamente significativa ( ver Gráfico 4.36b e Tabela II.18b, no Anexo IIB ).

% matched lemmas: 
\end{textsample}
